\rtmtight
\begin{tblr}{width=\textwidth,colspec={|Q[c,m,2.1cm]|X[l,m]|},hlines,vlines,hspan=minimal,rowsep=1.5pt,colsep=3pt,row{1}={bg=headgray,font=\bfseries}}
\SetCell{c,m}\hfil\logo\hfil & \textbf{POLITEKNIK NEGERI MALANG}\par\textbf{JURUSAN TEKNOLOGI INFORMASI}\par\textbf{PROGRAM STUDI: D4 TEKNIK INFORMATIKA} \\
\SetCell[c=2]{l} \textbf{RENCANA TUGAS MAHASISWA (RTM) DAN RUBRIK PENILAIAN} & \\
Nama Mata Kuliah & Rekayasa Perangkat Lunak \\
Kode Mata Kuliah & RTI252005 \quad \textbf{sks / Jam perminggu:} 2 SKS/4 jam \quad \textbf{Semester:} 2 \\
Dosen Pengampu & \begin{enumerate}\item Eka Larasati Amalia, S.ST., M.T.\item Rokhimatul Wakhidah, S.Pd., M.T.\item Vivi Nur Wijayaningrum, S.Kom., M.Kom.\item Candra Bella Vista, S.Kom., M.T.\item Dr. Eng. Banni Satria Andoko, S.Kom., M.MSI.\end{enumerate} \\
\end{tblr}
\assessmentblock{Tugas Dokumen SRS, Desain Sistem, dan Laporan Pengujian}{Tugas terstruktur (dokumen proyek)}{SCPMK0607-01402, SCPMK0211-01403, SCPMK0802-01401}{Mahasiswa menyusun dokumen spesifikasi kebutuhan perangkat lunak (SRS) hasil elisitasi, dokumen desain sistem (SDD) berbasis UML, dan laporan pengujian (Test Report) untuk proyek yang dikembangkan.}{Melakukan simulasi wawancara/elisitasi kebutuhan, mendokumentasikan SRS standar IEEE 830, mentransformasi SRS menjadi diagram UML dengan tools (StarUML/Draw.io), menyusun skenario uji, melakukan cross-testing, dan melaporkan bug.}{Dokumen SRS, dokumen desain sistem (SDD/UML), dan laporan pengujian (Test Report).}{Kelengkapan dan ketidakambiguan dokumen SRS & 10 & Kebutuhan fungsional dan non-fungsional teridentifikasi lengkap dari hasil elisitasi, dinyatakan tanpa ambigu, dan mengikuti struktur IEEE 830. & Kebutuhan utama terdokumentasi, tetapi sebagian pernyataan masih ambigu atau struktur dokumen belum lengkap. & Kebutuhan tidak lengkap, banyak pernyataan ambigu, atau dokumen tidak mengikuti standar. \\ \\
Ketepatan dan konsistensi desain UML (SDD) & 5 & Sintaks UML tepat; Use Case, Activity, Sequence, dan Class Diagram konsisten satu sama lain dan tertelusur ke SRS. & Diagram utama tersedia dan sebagian besar benar, tetapi terdapat inkonsistensi antar diagram atau dengan SRS. & Sintaks UML banyak yang salah, diagram tidak konsisten, atau tidak terkait dengan kebutuhan. \\ \\
Kelengkapan skenario uji dan laporan bug & 5 & Test Case mencakup skenario positif dan negatif, temuan bug akurat, dan laporan bug jelas serta dapat direproduksi. & Skenario uji tersedia tetapi cakupan terbatas, atau laporan bug kurang jelas. & Skenario uji tidak lengkap, temuan tidak akurat, atau laporan bug tidak dapat ditindaklanjuti. \\}

\assessmentblock{Kuis Konsep RPL, SDLC, dan Metodologi Pengembangan}{Kuis (tes objektif dan/atau esai singkat, closed book)}{SCPMK0802-01401}{Mahasiswa mengerjakan kuis materi pertemuan 1-4 yang mencakup konsep dasar RPL, tahapan SDLC, dan metodologi preskriptif (Waterfall, Prototyping, Spiral, RUP).}{Mengerjakan soal pilihan ganda dan/atau esai singkat secara mandiri, jujur, dan disiplin sesuai alokasi waktu ujian.}{Lembar jawaban kuis (daring/luring).}{Ketepatan menjelaskan tahapan SDLC & 1.5 & Tahapan SDLC dijelaskan lengkap, urut, dan disertai aktivitas serta luaran tiap tahapan. & Tahapan SDLC sebagian besar benar, tetapi urutan atau aktivitas tiap tahapan belum lengkap. & Tahapan SDLC keliru atau hanya menyebut istilah tanpa penjelasan. \\ \\
Ketepatan membedakan karakteristik metodologi & 1 & Karakteristik Waterfall, Prototyping, Spiral, dan RUP dibedakan tepat beserta konteks kapan masing-masing digunakan. & Sebagian karakteristik metodologi dijelaskan benar, tetapi konteks penggunaannya belum tepat. & Karakteristik metodologi tertukar atau jawaban tidak menunjukkan pemahaman. \\}

\assessmentblock{UTS Analisis dan Desain Sistem}{UTS (soal uraian dan studi kasus perancangan)}{SCPMK0802-01401, SCPMK0211-01403}{Mahasiswa mengerjakan soal studi kasus analisis dan desain sistem yang mencakup prinsip dasar RPL dan SDLC, metodologi preskriptif dan Agile, analisis kebutuhan (SRS), serta pemodelan visual sistem (UML).}{Mengerjakan soal uraian dan studi kasus perancangan secara mandiri (evaluasi sumatif) sesuai alokasi waktu ujian.}{Lembar jawaban UTS berisi analisis kasus dan rancangan diagram UML.}{Ketepatan pemilihan metodologi pengembangan & 2.5 & Alasan pemilihan metodologi (Agile vs preskriptif) dijelaskan berbasis karakteristik kasus, risiko, dan kejelasan kebutuhan. & Pemilihan metodologi benar, tetapi justifikasinya belum lengkap atau kurang terkait kasus. & Pemilihan metodologi keliru atau tanpa justifikasi yang relevan. \\ \\
Ketepatan transformasi kebutuhan menjadi diagram UML & 5 & Kebutuhan pada kasus ditransformasi menjadi Use Case, Activity, dan Class Diagram dengan notasi tepat dan konsisten antar diagram. & Diagram utama benar, tetapi terdapat kesalahan notasi atau inkonsistensi antar diagram. & Diagram tidak sesuai kebutuhan kasus atau notasi banyak yang salah. \\}

\assessmentblock{Case Method Analisis Model Pengembangan Perangkat Lunak}{Case Method (laporan kelompok)}{SCPMK0802-01401}{Mahasiswa menganalisis skenario proyek untuk membandingkan model Prototyping, Spiral, dan RUP serta menentukan metodologi yang paling sesuai dengan karakteristik kasus.}{Mengkaji literatur model SDLC, berdiskusi kelompok membandingkan model, menganalisis skenario proyek fiktif, dan menyusun dokumen analisis perbandingan beserta alasan pemilihan metode.}{Dokumen analisis studi kasus (laporan kelompok).}{Ketepatan menjelaskan karakteristik model pengembangan & 2 & Karakteristik Prototyping, Spiral, dan RUP dijelaskan tepat beserta kelebihan dan kekurangan masing-masing. & Karakteristik model sebagian besar benar, tetapi kelebihan/kekurangan belum lengkap. & Karakteristik model keliru atau tidak dibedakan antar model. \\ \\
Ketepatan analisis kasus dan pemilihan metodologi & 2 & Tingkat risiko dan kejelasan kebutuhan proyek teridentifikasi tepat sehingga metodologi terpilih sesuai karakteristik kasus. & Analisis kasus cukup tepat, tetapi keterkaitan risiko/kebutuhan dengan metodologi terpilih belum kuat. & Analisis kasus dangkal atau metodologi terpilih tidak sesuai kasus. \\ \\
Logika argumentasi pemilihan metode & 1 & Justifikasi pemilihan metode runtut, berbasis bukti dari kasus, dan mampu dipertahankan dalam diskusi. & Argumentasi cukup jelas, tetapi sebagian klaim tidak didukung bukti kasus. & Argumentasi tidak logis atau hanya mengulang teori tanpa kaitan ke kasus. \\}

\assessmentblock{PBL Proyek Pengembangan Perangkat Lunak}{Project Based Learning}{SCPMK0802-01401, SCPMK0602-01404, SCPMK0205-01405}{Mahasiswa mengembangkan produk perangkat lunak secara tim dalam siklus hidup proyek: inisiasi Scrum, perancangan arsitektur dan pola desain, manajemen alur kerja Kanban, dan implementasi kode.}{Membentuk tim Scrum dan menyusun Product Backlog, melakukan Sprint Planning, merancang arsitektur dan menerapkan Design Pattern, mengelola papan Kanban dengan WIP limits dan daily stand-up, serta melakukan coding kolaboratif dengan Git dan code review.}{Dokumen Product Backlog, project blueprint, project board (Trello/Jira/GitHub Projects), source code pada repository Git, dan demo aplikasi.}{Kelengkapan artefak Scrum dan Product Backlog & 10 & User Stories memenuhi kriteria INVEST, prioritas backlog jelas, peran tim sesuai, dan definisi Done (DoD) tegas. & Artefak Scrum tersedia, tetapi kualitas User Stories, prioritas, atau DoD belum konsisten. & Artefak Scrum tidak lengkap atau prosedur events Scrum tidak dijalankan. \\ \\
Kedisiplinan manajemen alur kerja Kanban & 5 & Kartu tugas terurai jelas, berpindah kolom sesuai status nyata (real-time), WIP limits dipatuhi, dan beban kerja tim merata. & Papan Kanban digunakan, tetapi pembaruan status atau pemerataan beban kerja belum konsisten. & Papan Kanban tidak mencerminkan pekerjaan nyata atau tidak digunakan oleh tim. \\ \\
Ketepatan arsitektur dan penerapan Design Pattern & 10 & Arsitektur (misal MVC) sesuai skala proyek, struktur proyek rapi, dan minimal satu Design Pattern diterapkan dengan benar dalam kode. & Arsitektur diterapkan, tetapi struktur proyek atau penerapan pola desain masih kurang tepat. & Arsitektur tidak jelas atau pola desain tidak diterapkan dalam kode. \\ \\
Kualitas implementasi kode & 20 & Kode modular sesuai standar Clean Code, aplikasi berjalan tanpa error/bug kritikal, dan kolaborasi Git serta code review terdokumentasi. & Aplikasi berjalan untuk fungsi utama, tetapi modularitas, kebersihan kode, atau jejak kolaborasi masih terbatas. & Aplikasi tidak berjalan sesuai fungsi, kode tidak modular, atau tidak ada bukti kolaborasi tim. \\}

\assessmentblock{UAS Presentasi dan Demo Produk Akhir}{UAS (project presentation dan demo produk)}{SCPMK0602-01404, SCPMK0802-01401}{Mahasiswa mempresentasikan hasil akhir produk perangkat lunak yang dikembangkan selama satu semester, mendemonstrasikan fitur, dan mempertahankan argumen teknis di hadapan penguji.}{Menyiapkan laporan akhir, mendemonstrasikan produk, menjelaskan penerapan arsitektur dan pola desain, menunjukkan bukti pengujian, dan menjawab pertanyaan penguji.}{Laporan akhir proyek, slide presentasi, demo aplikasi, serta dokumen maintenance (User Guide dan Tech Doc).}{Ketepatan penerapan arsitektur dan pola desain dalam kode & 10 & Arsitektur dan Design Pattern terlihat konsisten dalam kode program, dan mahasiswa mampu menjelaskan alasan teknisnya saat defense. & Arsitektur dan pola desain diterapkan, tetapi penjelasan teknis atau konsistensinya dalam kode masih lemah. & Arsitektur/pola desain tidak terlihat dalam kode atau tidak dapat dijelaskan. \\ \\
Kelengkapan skenario pengujian dan kualitas produk & 5 & Skenario pengujian lengkap, bukti pengujian tersedia, dan demo berjalan tanpa bug berarti. & Pengujian dilakukan, tetapi cakupan skenario terbatas atau masih muncul bug saat demo. & Tidak ada bukti pengujian atau demo gagal menunjukkan fungsi utama. \\ \\
Kelengkapan dokumen maintenance & 5 & User Guide dan Tech Doc lengkap, jelas, dan menunjukkan pemahaman siklus hidup software jangka panjang. & Dokumen maintenance tersedia, tetapi isi atau kejelasannya masih terbatas. & Dokumen maintenance tidak tersedia atau tidak dapat digunakan. \\}

\begin{tblr}{width=\textwidth,colspec={|Q[l,m,3.2cm]|X[l,m]|},hlines,vlines,hspan=minimal,row{1-3}={bg=headgray},rowsep=1.5pt,colsep=3pt}
Jadwal Pelaksanaan: & Minggu 1--16 sesuai RPS. Setiap evaluasi memiliki RTM dan rubrik multi-kriteria. \\
Lain-Lain Yang Diperlukan: & LMS, repository Git, tools pemodelan (StarUML/Draw.io), tools manajemen proyek (Trello/Jira/GitHub Projects), dan perangkat presentasi. \\
Pustaka: & Mengacu pustaka utama dan pendukung pada bagian RPS. \\
\end{tblr}
